\section{Thinking and Arguments}

\begin{quotex}
In those days it was more difficult to be reckoned as a thinker because there was no printing press, and it was not an easy thing to take part in such controversies as that between Nominalism and Realism. \emph{Anyone who ventured into this field had to be better prepared, according to the ideas of those times, than is required of people who engage in controversies nowadays.} An immense amount of penetration was necessary in order to plead the cause of Realism, and it was equally so with Nominalism.
\flright{\textsc{Rudolf Steiner}, \emph{Human and Cosmic Thought}}
\end{quotex}


Rudolf Steiner was describing the intellectual atmosphere in the Middle Ages. Books were rare and expensive, there was not a myriad of academic journals, and, of course, no conception of a blog which, had it even been anticipated, would have only aroused derision. Exalted thought was restricted to a very few, those who were intelligent, logical, and knowledgeable. Nevertheless, the debate between Nominalism and Realism ended in aporia, inconclusive. Like chess masters who know the game openings for a dozen or more moves, philosophical debaters know how to refute each other's arguments, so debates necessarily end in a draw.

Since no philosophical position can establish itself apodictically, the question arises (or should arise): ``On what basis does one adopt a world view?'' We get a clue from this confession in a recent blog post I bumped into:

\begin{quotex}
I try, without the aid of any particular theory or ideology, to understand things, and to form a consistent, non-contradictory view of the world.
\end{quotex}

In other words, by ignoring the greatest minds in history, this blogger is nevertheless able to form a worldview that is both logically non-contradictory and consistent with the observed facts. This is Kantianism once again, the disease of the modern mind; the rationale for this worldview is psychological, not logical. Now he is an intelligent and educated man, so he realizes he is fabricating a worldview. The common man, like him, also ignores the greatest minds, although he forms a worldview without much concern for consistency or non-contradiction.


In \textit{Human and Cosmic Thought}\footnote{\url{https://wn.rsarchive.org/Lectures/GA151/English/RSP1991/HCT991_index.html}}, Steiner identifies twelve possible worldviews which he associates with the signs of the Zodiac. They are:

\begin{itemize}


\item Ares: Idealism

\item Taurus: Rationalism

\item Gemini: Mathematism

\item Cancer: Materialism

\item Leo: Sensationalism

\item Virgo: Phenomenalism

\item Libra: Realism

\item Scorpio: Dynamism

\item Sagittarius: Monadism

\item Capricorn: Spiritism

\item Aquarius: Pneumatism

\item Pisces: Psychism
\end{itemize}


He points out that no one system of philosophy is able to eliminate the others, since each of them is plausible, at least within certain limits. Nevertheless, the game goes on, and people dispute endlessly. It is striking how a man will admit his inferiority in sports, since his abilities in comparison to the better players is plainly visible and cannot be ignored. However, then it comes to the most subtle and difficult philosophical issues, every man considers himself an infallible expert, while lacking, to a greater or lesser extent, the qualities of intelligence, rationality and knowledge. As Steiner puts it:

\begin{quotex}
Men dispute about world-outlooks, and why, on the other hand, they ought not to dispute but would do much better to understand why it happens that people have different world-outlooks.
\end{quotex}


The first thing to note, and few do, is that a ``thought'' is akin to a perception. An idea falls on the brain and becomes a thought. The twelve basic worldviews are analogous to constellations: related thoughts will tend to group together. As they bind themselves together, they will then tend to create a chain mail that will occlude other competing worldviews from entering our consciousness. That is why the first trial\footnote{\url{https://gornahoor.net/?p=1902}} is to transcend all our common patterns of thought.

Besides the twelve zodiacal worldviews, Steiner lists seven ``moods'' which correspond to planetary influences and affect the worldviews in different ways:

\begin{itemize}


\item Saturn: Gnosis

\item Jupiter: Logicism

\item Mars: Voluntarism

\item Sun: Empiricism

\item Venus: Mysticism

\item Mercury: Transcendentalism

\item Moon: Occultism
\end{itemize}


This scheme is part of the Hermetic and Rosicrucian Tradition, as can be seen in the Cosmic Hierarchy\footnote{\url{https://www.meditationsonthetarot.com/the-cosmic-hierarchy}} of Robert Fludd. The constellations, or worldviews, are represented by the \emph{Caelum Stellatum} layer, and the moods by the planets.

As the Cosmic Hierarchy shows, the real source of wordless thought is from the angelic hierarchies (and above). Insofar as the individual consciousness is disturbed by negative emotions, disordered desires, and intransigent worldviews, these higher thoughts will never take root in our consciousness. This process and the clarification of consciousness is described in The Word is Made Flesh\footnote{\url{https://www.meditationsonthetarot.com/the-word-is-made-flesh}}.

NOTE: Hermetism encompasses three Traditional sciences: alchemy, astrology, and theurgy (magic). In his The Hermetic Tradition, Julius Evola suggests that what he did for alchemy, needs to be done also for astrology and theurgy. This post is a possible beginning for whoever would like to take on that task.


\flrightit{Posted on 2011-05-20 by Cologero}

\begin{center}* * *\end{center}

\begin{footnotesize}
\begin{sffamily}
\texttt{William Bishop on 2012-06-26 at 09:44 said:}

Looking at the chart provided re. 12 paths.

It does not seem conducive that Gemini rule Math and Virgo Phenominalsim whereas Mercury is realted to Transcendentalism? Likewise Jupiter anethema in Gemini is ruling Logic and Sagitarius Monadism and Pisces Psychism. Is this supposed to be some paradoxical joke or something? Please someone explain if they think they can or even if it is worth while.

\hfill

\texttt{Cologero on 2012-06-26 at 18:48 said:}

William, the chart was taken from Rudolf Steiner's \textit{Human and Cosmic Thought}\footnote{\url{http://wn.rsarchive.org/Lectures/GA151/English/RSP1961/HuCoTh_index.html}}, where he offers justification for his schema. It may be fruitful to discuss his reasons.

Nevertheless, the main point of the article, especially as it relates to Hermetism, is that the traditional science of astrology is not a primitive precursor to the empirical science of astronomy. Rather, its real intent is to catalog different worldviews, that can bind us as long as we are unaware of the true source of the ideas we assume to be our ``own''.

\end{sffamily}
\end{footnotesize}

